\documentclass[11pt,a4paper]{article}
\usepackage[utf8]{inputenc}
\usepackage[vietnamese]{babel}
\usepackage{amsmath,amssymb}
\usepackage[version=4]{mhchem}
\usepackage{geometry}
\geometry{top=1.5cm,bottom=1.2cm,left=1.4cm,right=1.4cm,headheight=28pt,headsep=12pt,footskip=18pt}
\usepackage{graphicx}
\usepackage{xcolor}
\usepackage{tabularx}
\usepackage{array}
\usepackage{tikz}
\usepackage{multicol}
\usepackage{fancyhdr}
\usepackage{enumitem}

\definecolor{hfRed}{HTML}{D32F2F}
\definecolor{hfGreen}{HTML}{8BC390}
\definecolor{mainGreen}{HTML}{2E7D32}

\usepackage{ex_test}

\pagestyle{fancy}
\fancyhf{}
\renewcommand{\headrulewidth}{0.75pt}
\renewcommand{\headrule}{\hbox to\headwidth{\color{hfGreen}\leaders\hrule height \headrulewidth\hfill}}
\renewcommand{\footrulewidth}{0.75pt}
\renewcommand{\footrule}{\hbox to\headwidth{\color{hfGreen}\leaders\hrule height \footrulewidth\hfill}}

\fancyhead[L]{\small\itshape\color{hfRed}Tài liệu lưu hành nội bộ}
\fancyhead[R]{\small\color{hfRed}Thầy \textbf{TRẦN VĂN THẠNH} - 0777.470.803}

\fancyfoot[L]{\footnotesize\itshape\color{hfRed}CS1: 6/15 Nguyễn Hoàng, P. Kim Long, TP Huế \quad|\quad CS2: 24 Đặng Thái Thân, TP Huế}
\fancyfoot[R]{\small\textbf{\color{hfRed}Trang \thepage}}

\setlength{\parindent}{0pt}
\setlength{\parskip}{2.0pt}

\def\loigiai#1{}

\begin{document}

\Opensolutionfile{ans}[ans-de10-hs]

\noindent
\begin{minipage}[t]{0.48\textwidth}
	\centering
	\textbf{SỞ GIÁO DỤC VÀ ĐÀO TẠO THỪA THIÊN HUẾ}\\
	\textbf{TRƯỜNG THPT HAI BÀ TRƯNG}\\
	\textbf{ĐỀ CHÍNH THỨC}\\
	\textit{(Đề thi có 04 trang)}
\end{minipage}\hfill
\begin{minipage}[t]{0.48\textwidth}
	\centering
	\textbf{\color{mainGreen}ĐỀ KIỂM TRA ĐỊNH KỲ CHƯƠNG 4 (POLYMER)}\\
	\textbf{MÔN: HOÁ HỌC -- LỚP 12}\\
	\textit{Thời gian làm bài: 50 phút (không kể thời gian phát đề)}\\
	\textbf{Mã đề thi: 110}
\end{minipage}


\vspace{0.1cm}
\noindent\textit{Họ, tên thí sinh: \dots\dots\dots\dots\dots\dots\dots\dots\dots\dots\dots\dots\dots\dots\dots\dots\dots\dots\dots\dots\dots\dots Lớp: \dots\dots\dots\dots SBD: \dots\dots\dots\dots}\\
\noindent\textit{(Cho biết nguyên tử khối: H = 1; C = 12; N = 14; O = 16; Na = 23; Al = 27; S = 32; Cl = 35,5; K = 39; Ca = 40; Fe = 56; Cu = 64; Zn = 65; Br = 80; Ag = 108; Ba = 137).}

\vspace{0.15cm}
\noindent
\textbf{PHẦN I (4,5 điểm -- 18 câu): Câu trắc nghiệm nhiều phương án lựa chọn.}\\
\textit{Thí sinh trả lời từ câu 1 đến câu 18. Mỗi câu hỏi thí sinh chỉ chọn một phương án.}

\setcounter{ex}{0}

\begin{ex}
Polymer nào sau đây thuộc loại polymer thiên nhiên?
\choice
{Tơ visco}
{Poly(vinyl chloride)}
{Tinh bột}
{Polyethylene}
\loigiai{}
\end{ex}

\begin{ex}
Polymer nào sau đây có công thức \ce{[-CH2-CH(CN)-]_n}?
\choice
{Poly(methyl methacrylate)}
{Polyethylene}
{Polyacrylonitrile}
{Poly(vinyl chloride)}
\loigiai{}
\end{ex}

\begin{ex}
Poly(vinyl chloride) được điều chế trực tiếp từ monomer nào sau đây?
\choice
{\ce{CH2=CH-CN}}
{\ce{CH2=CH2}}
{\ce{CH2=CH-CH=CH2}}
{\ce{CH2=CH-Cl}}
\loigiai{}
\end{ex}

\begin{ex}
Polyacrylonitrile được điều chế trực tiếp từ monomer nào sau đây?
\choice
{\ce{CH2=CH-Cl}}
{\ce{CH2=CH2}}
{\ce{CH2=CH-CN}}
{\ce{CH2=CH-CH=CH2}}
\loigiai{}
\end{ex}

\begin{ex}
(SBT -- KNTT). Polypropylene là chất dẻo được sử dụng phổ biến thứ hai sau polyethylene. Trùng hợp chất nào sau đây thu được polypropylene?
\choice
{\ce{CH2=CH-Cl}}
{\ce{CH2=CH2}}
{\ce{CH2=CH-C6H5}}
{\ce{CH2=CH-CH3}}
\loigiai{}
\end{ex}

\begin{ex}
(SBT -- KNTT). Trùng hợp monomer \ce{CH2=CH-Cl} thu được chất dẻo nào sau đây?
\choice
{PE}
{PP}
{PVC}
{PS}
\loigiai{}
\end{ex}

\begin{ex}
(SBT -- KNTT). Tơ sợi nào sau đây thuộc loại tơ tự nhiên?
\choice
{Sợi bông}
{Nitron}
{Nylon-6,6}
{Cellulose acetate}
\loigiai{}
\end{ex}

\begin{ex}
(SBT -- KNTT). Trùng hợp chất nào sau đây thu được polyacrylonitrile dùng để sản xuất tơ nitron?
\choice
{\ce{CH2=CH-Cl}}
{\ce{CH2=CH-CN}}
{\ce{CH2=CH2}}
{\ce{CH2=CH-CH3}}
\loigiai{}
\end{ex}

\begin{ex}
(SBT -- KNTT). Trùng hợp chất nào sau đây thu được cao su isoprene?
\choice
{\ce{CH2=CH-CH=CH2}}
{\ce{CH2=CCl-CH=CH2}}
{\ce{CH2=C(CH3)-CH=CH2}}
{\ce{CH2=C(CH3)-CCl=CH2}}
\loigiai{}
\end{ex}

\begin{ex}
Dãy polymer đều được điều chế từ phản ứng trùng ngưng là
\choice
{Teflon, polyethylene, PVC}
{Cao su buna, nylon-7, tơ acetate}
{Nylon-6, poly(vinyl alcohol), thủy tinh plexiglass}
{Nhựa PPF, nylon-7, tơ lapsan}
\loigiai{}
\end{ex}

\newpage

\begin{ex}
Trong các polymer: polyethylene, poly(methyl methacrylate), tinh bột, tơ tằm, tơ nylon-6, số polymer thiên nhiên là
\choice
{4}
{2}
{5}
{3}
\loigiai{}
\end{ex}

\begin{ex}
Cho các polymer sau: \ce{[-CH2-CH2-]_n}; \ce{[-CH2-CH=CH-CH2-]_n}; \ce{[-NH-CH2-CO-]_n}. Công thức của các monomer để khi trùng hợp hoặc trùng ngưng tạo ra các polymer trên lần lượt là
\choice
{\ce{CH2=CHCl}, \ce{CH3-CH=CH-CH3}, \ce{CH3-CH(NH2)-COOH}}
{\ce{CH2=CH2}, \ce{CH2=CH-CH=CH2}, \ce{H2N-CH2-COOH}}
{\ce{CH2=CH2}, \ce{CH3-CH=C=CH2}, \ce{NH2-CH2-COOH}}
{\ce{CH2=CH2}, \ce{CH3-CH=CH-CH3}, \ce{H2N-CH2-CH2-COOH}}
\loigiai{}
\end{ex}

\begin{ex}
Dãy nào sau đây chỉ gồm các polymer tổng hợp?
\choice
{Polypropylene, cellulose, nylon-7, nylon-6,6}
{Polypropylene, polybutadiene, nylon-7, nylon-6,6}
{Polypropylene, tinh bột, nylon-7, cao su thiên nhiên}
{Tinh bột, cellulose, cao su thiên nhiên, polybutadiene}
\loigiai{}
\end{ex}

\begin{ex}
Cho sơ đồ chuyển hoá:
\centerline{$\text{Glucose} \xrightarrow{\text{enzyme}} \text{X} \xrightarrow{\text{ZnO},\, 450^\circ\text{C}} \text{Y} \xrightarrow{xt,\, t^\circ,\, p} \text{Cao su Buna.}$}
Hai chất X, Y lần lượt là
\choice
{\ce{CH3CH2OH} và \ce{CH3CHO}}
{\ce{CH3CH2OH} và \ce{CH2=CH2}}
{\ce{CH2CH2OH} và \ce{CH3-CH=CH-CH3}}
{\ce{CH3CH2OH} và \ce{CH2=CH-CH=CH2}}
\loigiai{}
\end{ex}

\begin{ex}
Nhóm các polymer được dùng làm tơ là
\choice
{Polyacrylonitrile, poly(methyl methacrylate)}
{Poly(vinyl chloride), polybuta-1,3-diene}
{Polyacrylonitrile, poly(hexamethylene adipamide)}
{Poly(hexamethylene adipamide), poly(vinyl chloride)}
\loigiai{}
\end{ex}

\begin{ex}
Nhận xét về tính chất vật lí chung của polymer nào dưới đây \textbf{không} đúng?
\choice
{Hầu hết là những chất rắn, không bay hơi, không có nhiệt độ nóng chảy xác định}
{Khi nóng chảy, đa số polymer cho chất lỏng nhớt, để nguội sẽ rắn lại gọi là chất nhiệt dẻo}
{Một số polymer không nóng chảy khi đun mà bị phân hủy, gọi là chất nhiệt rắn}
{Polymer không tan trong nước và trong bất kỳ dung môi nào}
\loigiai{}
\end{ex}

\begin{ex}
Cho các phát biểu sau:\\
(1) Tơ viscose, tơ acetate đều thuộc loại tơ tổng hợp.\\
(2) Polyethylene và poly(vinyl chloride) là sản phẩm của phản ứng trùng hợp.\\
(3) Tơ nylon-6,6 được điều chế từ hexamethylenediamine và stearic acid.\\
(4) Cao su thiên nhiên không tan trong nước cũng như trong xăng, benzene.\\
(5) Tơ nitron (olon) được tổng hợp từ vinyl cyanide (acrylonitrile).\\
Số phát biểu đúng là
\choice
{4}
{3}
{5}
{2}
\loigiai{}
\end{ex}

\begin{ex}
Một polymer X được xác định có phân tử khối là 39062,5 amu với hệ số trùng hợp để tạo nên polymer này là 625. Polymer X là
\choice
{Poly(vinyl chloride) (PVC)}
{Polypropylene (PP)}
{Polyethylene (PE)}
{Polystyrene (PS)}
\loigiai{}
\end{ex}

\newpage

\noindent
\textbf{PHẦN II (4,0 điểm -- 4 câu): Câu trắc nghiệm đúng sai.}\\
\textit{Thí sinh trả lời từ câu 1 đến câu 4. Trong mỗi ý a), b), c), d) ở mỗi câu, thí sinh chọn đúng hoặc sai.}

\setcounter{ex}{0}

\begin{ex}
Hộp xốp đựng thực phẩm chế biến sẵn thường làm bằng polystyrene (PS).
\choiceTF
{Polystyrene được điều chế bằng phản ứng trùng ngưng}
{Có thể sử dụng hộp nhựa PS để đựng thức ăn nóng hoặc cho vào lò vi sóng}
{Styrene dùng để điều chế ra PS}
{Khi đồng trùng hợp buta-1,3-diene với styrene ta thu được cao su buna-S}
\loigiai{}
\end{ex}

\begin{ex}
(SBT -- CĐ). Tơ sợi là một nguyên liệu quan trọng trong ngành dệt may và sản xuất vật liệu. Tơ sợi được sản xuất từ các nguồn nguyên liệu tự nhiên hoặc tổng hợp. Các loại tơ sợi phổ biến bao gồm tơ sợi tự nhiên như tơ tằm, lông cừu, sợi cotton và tơ sợi tổng hợp như nylon,\dots{} Những phát biểu nào sau đây là đúng?
\choiceTF
{Tơ nylon-6,6 và tơ capron thuộc loại tơ polyamide}
{Tơ nylon, tơ tằm, tơ visco đều bền với nhiệt độ}
{Quần áo được dệt bằng sợi len lông cừu, tơ tằm không nên giặt với xà phòng có độ kiềm cao}
{Tơ capron và tơ olon đều có thành phần chứa nhóm \ce{-CO-NH-}}
\loigiai{}
\end{ex}

\begin{ex}
Poly(ethylene terephthalate) (viết tắt là PET) là một polymer được điều chế từ terephthalic acid và ethylene glycol. PET được sử dụng để sản xuất tơ, chai đựng nước uống, hộp đựng thực phẩm. Để thuận lợi cho việc nhận biết, sử dụng và tái chế thì các đồ nhựa làm từ vật liệu chứa PET thường được in kí hiệu như hình bên.
\begin{center}
	\includegraphics[width=2.4cm]{images_crop_de10/p2_cau3_pet.png}
\end{center}
\choiceTF
{PET thuộc loại polyester}
{Tơ được chế tạo từ PET thuộc loại tơ tổng hợp}
{Trong một mắt xích PET, phần trăm khối lượng carbon là 62,5\%}
{Phản ứng tổng hợp PET từ terephthalic acid và ethylene glycol thuộc loại phản ứng trùng hợp}
\loigiai{}
\end{ex}

\begin{ex}
Công ty The Goodyear Tire \& Rubber là một trong những công ty lốp xe lớn nhất thế giới khởi lập năm 1898. Năm 1971, lốp Goodyear trở thành bánh xe đầu tiên lăn trên Mặt Trăng\dots{} Tên công ty được đặt theo tên của nhà tiên phong Charles Goodyear, người khám phá ra phương pháp kết hợp giữa nguyên tố S (lưu huỳnh) với cao su để tạo ra một loại cao su mới.
\choiceTF
{Loại cao su này có tên là cao su lưu hóa}
{Cao su lưu hóa có cấu trúc mạch phân nhánh}
{Cao su lưu hóa làm tăng cao tính bền hơn với nhiệt và các tác nhân khác, độ đàn hồi cao, chống thấm khí, chống thấm ẩm, tăng khả năng chịu được sự ma sát, va chạm tốt hơn cao su thiên nhiên}
{Cao su lưu hóa làm tăng tính bền cơ học nhờ các cầu nối trisulfide (\ce{-S-S-S-}) giữa các phân tử polyisoprene}
\loigiai{}
\end{ex}

\newpage

\noindent
\textbf{PHẦN III (1,5 điểm -- 6 câu): Câu trắc nghiệm yêu cầu trả lời ngắn.}\\
\textit{Thí sinh trả lời từ câu 1 đến câu 6.}

\setcounter{ex}{0}

\begin{ex}
(SBT -- CTST). Có bao nhiêu loại polymer có thể tạo thành từ adipic acid, 1,2-ethylene diamine và ethylene glycol bằng phản ứng trùng ngưng?
\par\noindent\textit{Trả lời:} \dotfill
\loigiai{}
\end{ex}

\begin{ex}
Cho dãy các polymer sau: polyethylene, cellulose, nylon-6,6, amylose, nylon-6, tơ nitron, polybuta-1,3-diene, tơ visco, poly(phenol formaldehyde), poly(vinyl acetate). Có bao nhiêu polymer tổng hợp trong dãy trên?
\par\noindent\textit{Trả lời:} \dotfill
\loigiai{}
\end{ex}

\begin{ex}
Khối lượng phân tử của một loại thủy tinh hữu cơ plexiglass là 25000 amu. Số mắt xích trong phân tử thủy tinh hữu cơ đó là bao nhiêu?
\par\noindent\textit{Trả lời:} \dotfill
\loigiai{}
\end{ex}

\begin{ex}
Khối lượng phân tử của một đoạn polyethylene (PE) là 28000 amu. Số mắt xích của đoạn polymer này là bao nhiêu?
\par\noindent\textit{Trả lời:} \dotfill
\loigiai{}
\end{ex}

\begin{ex}
(SBT -- CTST). Cellulose triacetate được điều chế từ cellulose và acetic anhydride. Tính khối lượng cellulose triacetate (theo đơn vị tấn) tạo thành từ 1 tấn cellulose nếu hiệu suất của quá trình điều chế là 67\%. (Làm tròn kết quả đến chữ số thập phân thứ hai).
\par\noindent\textit{Trả lời:} \dotfill
\loigiai{}
\end{ex}

\begin{ex}
Một loại cao su Buna-S có phần trăm khối lượng carbon là 90,225\%; $m$ gam cao su này cộng tối đa với 9,6 gam bromine. Giá trị của $m$ (theo đơn vị gam) là bao nhiêu? (Làm tròn kết quả đến chữ số thập phân thứ hai).
\par\noindent\textit{Trả lời:} \dotfill
\loigiai{}
\end{ex}


\Closesolutionfile{ans}

\end{document}
